% !TEX TS-program = xelatex

\documentclass[
    10pt,
    a5paper,
    oneside,
    openany
]{book}

% -------------------------------------------------
% Formatting
% -------------------------------------------------

% \usepackage[blue]{formatting}
\usepackage[blue]{formatting}
\usepackage{mhchem}

% -------------------------------------------------
% Book information
% -------------------------------------------------

\newcommand{\booktitle}{Medicine}
\newcommand{\booksubtitle}{at Swansea University}
\newcommand{\bookauthor}{Jack Beda}
\newcommand{\bookinstitution}{}
\newcommand{\bookedition}{Compiled \today}

% \newcommand{\centredurl}[1]{%
%     \par\smallskip
%     {\centering\url{#1}\par}%
%     \smallskip
% }



\hypersetup{
    pdftitle={\booktitle},
    pdfauthor={\bookauthor},
    pdfsubject={Medicine course notes},
    pdfkeywords={medicine, medical notes, revision, handbook}
}

% LaTeX will look in these folders when loading images.
\graphicspath{
    {images/}
    {chapters/images/}
}

\setcounter{tocdepth}{2}
\setcounter{secnumdepth}{3}

\begin{document}

% -------------------------------------------------
% Front matter
% -------------------------------------------------

\frontmatter

\begin{titlepage}
    \thispagestyle{empty}

    \begin{tikzpicture}[remember picture,overlay]

        % Full-page primary-colour background.
        \fill[PrimaryColor]
            (current page.south west)
            rectangle
            (current page.north east);

        % Decorative accent circles.
        \fill[AccentColor,opacity=0.28]
            ([xshift=-12mm,yshift=-18mm]current page.north east)
            circle[radius=40mm];

        \fill[SoftAccentColor,opacity=0.15]
            ([xshift=5mm,yshift=4mm]current page.south west)
            circle[radius=31mm];

        % Subtle horizontal divider.
        \draw[
            white,
            opacity=0.14,
            line width=1.2pt
        ]
            ([xshift=16mm,yshift=25mm]current page.south west)
            --
            ([xshift=-16mm,yshift=25mm]current page.south east);

        % Accent-colour band along the bottom.
        \fill[AccentColor]
            (current page.south west)
            rectangle
            ([yshift=7mm]current page.south east);

    \end{tikzpicture}

    \color{white}

    \vspace*{20mm}

    \begin{flushleft}

        \vspace{8mm}

        {\sffamily\bfseries\fontsize{31}{35}\selectfont
            \booktitle
            \par
        }

        \vspace{5mm}

        {\sffamily\Large\color{white!82}
            \booksubtitle
            \par
        }

        \vspace{8mm}

        {\color{AccentColor}\rule{36mm}{1.5pt}}

    \end{flushleft}

    \vfill

    \begin{flushleft}

        {\sffamily\large\bfseries
            \bookauthor
            \par
        }

        \vspace{2mm}

        {\sffamily\normalsize\color{white!80}
            \bookinstitution
            \par
        }

        \vspace{1mm}

        {\sffamily\normalsize\color{white!80}
            \bookedition
            \par
        }

        \vspace{12mm}

    \end{flushleft}

    \vspace{6mm}

\end{titlepage}

\normalcolor


\pagestyle{notes}
\clearpage

\tableofcontents

\clearpage
\chapter{Introduction}

This document contains my notes from Swansea University's graduate-entry
medicine course. The notes are accessible through GitHub:

\centredurl{https://github.com/jfbeda/medicine}

You are encouraged to edit and contribute to the notes as you see fit.
My contact details are available on my website:

\centredurl{https://beda.ca/jack}

\section{Organisation}

The main body of the handbook is divided into chapters. Each chapter is
stored as a separate file inside the \texttt{chapters} folder.

For example:

\begin{itemize}
    \item \texttt{chapters/02-cardiovascular.tex}
    \item \texttt{chapters/03-respiratory.tex}
    \item \texttt{chapters/04-gastrointestinal.tex}
\end{itemize}

The chapters are imported in the required order from
\texttt{main.tex}.


% -------------------------------------------------
% Main chapters
% -------------------------------------------------

\mainmatter
\pagestyle{notes}

% Add chapter files below in the order in which they
% should appear in the book. Do not include ".tex".

\input{chapters/mandatory_training}
\chapter{Resources and Useful Links}
\label{ch:useful_links}

\section{Administration links}
\label{sec:administration_links}

\begin{itemize}
    \item \href{https://scs.swan.ac.uk/gem/UnitLogIn.php}{\textbf{Swansea GEM portal}}
    \begin{itemize}
    \item \href
        {https://scs.swan.ac.uk/gem/feedback/students/feedback_index.php}
        {\textbf{Student feedback}}
    \item \href
        {https://scs.swan.ac.uk/gem/isr/students/isr_loa.php}
        {\textbf{Leave of absence request}}
    \end{itemize}
    \item \href{https://swansea-uk.libcal.com/reserve/splgroupstudy}{\textbf{Library group study room booking}}
    \item \href{https://www.swansea.ac.uk/finance-swansea-university/paying-tuition-fees-and-other-information/pay-fees-international-students/}{\textbf{International guidance on paying tuition fees}}
    \item \href{https://intranet.swan.ac.uk/studentprofile/profilesummarydetails.aspx}{\textbf{Intranet Student profile}}
\end{itemize}

\section{Learning resources}
\label{sec:learning_resources}
\begin{itemize}
    \item \href{https://3d4medical.com/}{\textbf{Complete anatomy} (activation code: 83548XFMGSQJ)}
    \item \href{https://dissectibleme.podbean.com/}{\textbf{Dissectable Me} (5 minute anatomy podcasts)}
    \item \href{https://www.kenhub.com/}{\textbf{Kenhub} (anatomy learning platform)}
    \item \href{https://teachmeanatomy.info/}{\textbf{Teach Me Anatomy} (anatomy learning platform)}
    \item \href{https://rhodridavies7.github.io/swansea-gem-anatomy/spotter.html}{\textbf{Swansea GEM anatomy spotter quiz}}
\end{itemize}


% \begin{itemize}
%     \item Moore's clinically oriented anatomy
%     \item Gray's anatomy for students
%     \item RHiMME track: \href{https://forms.cloud.microsoft/Pages/ResponsePage.aspx?id=LrXKu76f1kOi859mxD3yaEv50t8V6HJGhLbdbv-llK9UMlVKTTlQMkNIS01OVFNaVDRUNVBYTlJOTi4u&origin=QRCodehttps://forms.cloud.microsoft/Pages/ResponsePage.aspx?id=LrXKu76f1kOi859mxD3yaEv50t8V6HJGhLbdbv-llK9UMlVKTTlQMkNIS01OVFNaVDRUNVBYTlJOTi4u&origin=QRCode}{click here} Due September 25, 5pm
% \end{itemize}

\chapter{Useful Contacts}

\begin{itemize}
    \item \href{mailto:Kenneth.mckeegan@swansea.ac.uk}{\textbf{Kenny McKeegan}}
    \item \href{mailto:v.c.rowe@swansea.ac.uk}{\textbf{Carl Rowe}}
        \begin{itemize}
            \item Academic Head of GEM / Programme Director
            \item \textbf{Office}: SUSIM building?
        \end{itemize}
    \item \href{mailto:Christian.Cobbold@Swansea.ac.uk}{\textbf{Christian Cobbold}}
        \begin{itemize}
            \item \textbf{Office}: Grove building, room 243, near the GEM common room.
        \end{itemize}
  \item \href{mailto:b.w.mason@swansea.ac.uk}{\textbf{Brendan Mason}}
    \begin{itemize}
    \item Head of literature review
      \item Happy to be emailed again if you haven't received a reply after \SI{48}{\hour}
    \end{itemize}
    \item \href{mailto:k.krishnamoorthy@swansea.ac.uk}{\textbf{Karpagam Krishnamoorthy}}
    \begin{itemize}
        \item  \textbf{Office}: Grove 154A
    \end{itemize}
\end{itemize}

\chapter{Anatomy}
\label{ch:anatomy}


\section{Anatomical terms}

\subsection{Anatomical planes}

\begin{itemize}
    \item Coronal (frontal) plane: divides the body into anterior and posterior halves.
    \item Sagittal plane: divides the body into left and right halves. The \term{median plane} divides the body into equal left and right halves, while a \term{parasagittal plane} divides the body into unequal left and right halves.
    \item Transverse (horizontal or axial) plane: divides the body into superior and inferior halves.
\end{itemize}

\subsection{Anatomical positional terms}

\begin{itemize}
    \item Superior
    \item Inferior 
    \item Cephalic (towards the head)\footnote{The terms \term{cephalic} and \term{caudal} are often used in embryology, where the head and tail of the embryo are more distinct than in adults.}
    \item Caudal (towards the tail)\footnotemark[\value{footnote}]
    \item Proximal (closer to the point of attachment)
    \item Distal (further from the point of attachment)
    \item Medial
    \item Lateral
    \item Posterior/dorsal
    \item Anterior/ventral
\end{itemize}

\subsection{Anatomical movement terms}

\begin{itemize}
    \item Supination: rotation of the forearm so that the palm faces anteriorly (or superiorly when the arm is outstretched)
    \item Pronation: rotation of the forearm so that the palm faces posteriorly (or inferiorly when the arm is outstretched)
    \item Inversion: rotation of the foot so that the sole faces medially
    \item Eversion: rotation of the foot so that the sole faces laterally
    \item Flexion: movement that decreases the angle between two body parts\footnote{\label{fn:angle} the `angle' being referred to here is the angle in the `natural' movement (e.g. at the hip it is the raising of the leg that decreases the angle on the anterior side, which is the side that is usually considered when describing flexion and extension).}
    \item Extension: movement that increases the angle between two body parts\footnotemark[\value{footnote}]
    \item Abduction: movement away from the midline of the body
    \item Adduction: movement towards the midline of the body
    \item Circumduction: circular movement of a body part, combining flexion, extension, abduction, and adduction
    \item Rotation: movement around a longitudinal axis, either medially (internal rotation) or laterally (external rotation)
    \item Elevation: movement of a body part superiorly
    \item Depression: movement of a body part inferiorly
    \item Protraction: movement of a body part anteriorly
    \item Retraction: movement of a body part posteriorly
    \item Opposition: movement of the thumb to touch the tips of the other fingers
    \item Reposition: movement of the thumb back to its anatomical position
    \item Dorsiflexion: movement of the foot so that the toes point superiorly
    \item Plantar flexion: movement of the foot so that the toes point inferiorly
    \item Palmar flexion: movement of the hand so that the palm faces inferiorly
\end{itemize}

\chapter{Diarrhea}
\label{ch:diarrhea}

\section{The Small Intestine}
Based on the video \textit{\href{https://www.youtube.com/watch?v=6BPMo8CbLCw&list=PLX2ZvPJ5dxu4}{The small intestine: How does the small intestine absorb so much?}} by Christian Cobbold.

\begin{itemize}
    \item Small intestine has many villi and microvilli to increase surface area for absorption.
    \item Each epithelial cell contains order \SI{3000} microvilli.
    \item Glucose is absorbed into cells via sodium-glucose transporters (SGLT1) that work by pumping sodium ions out of the cell (towards the interstitial fluid), creating a sodium gradient (with respect to the intestinal fluid), that results in sodium uptake from the intestinal fluid where glucose is co-transported with sodium into the cell.
    \item Fat globules are emulsified by bile salts which are then (somehow?) turned into micelles (whatever they happen to be?) which are then turned into water-soluble chylomicrons (diameter order \SI{1}{\micro \meter}) that are absorbed into the lymphatic system.
\end{itemize}

\section{The Large Intestine}
Based on the video \textit{\href{https://www.youtube.com/watch?v=aRas2Rv5Ssc&list=PLX2ZvPJ5dxu4}{Large intestine: Absorption and defaecation}} by Christian Cobbold.

\begin{itemize}
    \item Age results in decreased gastrointestinal motility, which can result in constipation.
    \item Age results in decreased large intestine muscle tone and thus straining when defecating, which can result in haemorrhoids.
    \item Age results in weakening sphincter muscles which results in acid reflux and heartburn.
\end{itemize}

\section{Fluid Movement Between Body Compartments}

Based on the video \textit{\href{https://www.youtube.com/watch?v=aLQci5SMXiQ&list=PLNIQOIcOf3y0}{Fluid movement between body compartments: Why doesn't all your fluid stay in the bloodstream?}} by Christian Cobbold.

\section{Oedema}
Based on the video \textit{\href{https://www.youtube.com/watch?v=Ya15ykZo6NA&list=PLW6XikaOqNG0}{Oedema: Why fluid accumulates in tissues }} by Christian Cobbold.


\chapter{Conditions}

\section{Clostridium difficile infection (c. diff)}
\label{cond:c-diff}

\section{Peptic ulcer}
\label{cond:peptic_ulcer}

Begins with a break in the epithelium which can result from e.g. \hyperref[cond:h-pylori]{h. pylori}, \hyperref[drugclass:NSAIDs]{NSAIDs}

30\% of ulcers are gastric, and 70\& are duedenal \cite[7:32]{CobboldPepticUlcerDisease}.

\section{Helicobacter pylori (h. pylori)}
\label{cond:h-pylori}

Often acquired through oral-oral or oral-fecal in children \cite{CobboldPepticUlcerDisease}

h. pylori can active G cells, which cause the release of more gastrin, and then more gastric acid, which results in mucosal injury \cite[t. 5:50]{CobboldPepticUlcerDisease}. 

Gastrin also causes a release of pepsinogen, which turns to pepsin under the influence of acid, which then digests exposed mucosal proteins \cite[t. 6:47]{CobboldPepticUlcerDisease}. Takeaway message is that pepsin contributes to ulcer deepening, not the cause. 

\section{Gastro-oesophageal reflux disease (GORD/GERD)}

GORD is any continuous or regular reflux. Anything that disrupts the lower-oesophageal sphincter can lead to GORD

Long term GORD can cause reflux oesophagitus (mucosal redness and itching). Eventually you can get bleeding and scarring in the oseophagus, and then extremely you can get oesophageal narrowing due to scarring and finally Barrett's oesophagus. \cite[9:--]{CobboldGastroOesophagealRefluxDisease}

\textbf{Causes}

\begin{itemize}
    \item Alcohol and caffeine can cause lower-oesophageal relaxation
\end{itemize}
\chapter{Drugs}

\section{Antibiotics}

\subsection{Penicillin}

\subsection{Amoxicillin}


\section{Proton Pump Inhibitors (PPIs)}

Examples include

\begin{itemize}
    \item Omeprazole
    \item Lansoprazole
    \item Esomeprazole
    \item Pantoprazole
\end{itemize}

PPIs travel through the stomach, and are absorbed in the small intestine into the blood. They act on gastric parietal cells to inhibit the release of gastric acid.

\textbf{Side effects} include:

\begin{itemize}
    \item Increased risk of \hyperref[condition:c-diff]{\textit{c. diff}} infection
\end{itemize}

\section{Alginates}

Examples include

\begin{itemize}
    \item Gaviscon
    \item Peptac
    \item Wind-eze
\end{itemize}

Forms a 'raft' over the stomach contents.

\section{Antacids}

\textbf{Mechanism of action}: neutralizes stomach acid with a chemical reaction.

\textbf{Examples include}:

\begin{itemize}
    \item Calcium carbonate (\ce{CaCO3})
    \item Magnesium carbonate (\ce{MgCO3})
\end{itemize}

For calcium carbonate, the reaction is as follows:

\begin{equation}
    \ce{CaCO3 + 2HCl -> CaCl2 + H2O + CO2}
\end{equation}

\section{Prokinetics}

\textbf{Examples include}:

\begin{itemize}
    \item Metoclopramide
    \item Domperidone
    \item Enythromycin (off-license)
\end{itemize}

\section{H2-receptor antagonists (H2RAs)}

\textbf{Mechanism of action}: binds to H2 receptors on gastric parietal cells, preventing histamine from binding and thus reducing gastric acid secretion.

\textbf{Examples include}:

\begin{itemize}
    \item Famotidine
    \item Cimetidine
\end{itemize}

\section{Non-Steroidal Anti-Inflammatory Drugs (NSAIDs)}

\textbf{Examples include}:

\begin{itemize}
    \item Ibuprofen
    \item Naproxen
    \item Diclofenac
\end{itemize}

\textbf{Mechanism of action}

Prostaglandins produce pain/inflammation/fever. If we inhibit COX-2, we can avoid producing prostaglandins, but as a side effect that prevents the protection of the gastric mucosa.
\chapter{Acronyms}

\begin{itemize}
    \item NGT: Nasogastric Tube (2026-09-11 lecture with Pramodh V)
    \item ORS: Oral Rehydration Solution (2026-09-11 lecture with Pramodh V)
    \item LOCS: Learning Opportunities in Clinical Settings
    \item CC: Cubic Centimeter
    \item SAM: Severe Acute Malnutrition (2026-09-11 lecture with Pramodh V)
    \item MUAC: Mid Upper Arm Circumference (2026-09-11 lecture with Pramodh V) (9.8 cm)
    \item SCM: Sternocleidomastoid %https://youtu.be/Qgk2K_OYxQE?si=CxnfleEz1k6SVLTn&t=10
    \item PC: Presenting Complaint
    \item HPC: History of presenting complaint
\end{itemize}


% Examples:
% \input{chapters/02-cardiovascular}
% \input{chapters/03-respiratory}
% \input{chapters/04-gastrointestinal}
% \input{chapters/05-neurology}
% \input{chapters/06-endocrinology}
% \input{chapters/07-renal}
% \input{chapters/08-infectious-diseases}
% \input{chapters/09-pharmacology}

% Optional organisation using parts:
%
% \part{Foundations}
% \chapter{Anatomy}
\label{ch:anatomy}


\section{Anatomical terms}

\subsection{Anatomical planes}

\begin{itemize}
    \item Coronal (frontal) plane: divides the body into anterior and posterior halves.
    \item Sagittal plane: divides the body into left and right halves. The \term{median plane} divides the body into equal left and right halves, while a \term{parasagittal plane} divides the body into unequal left and right halves.
    \item Transverse (horizontal or axial) plane: divides the body into superior and inferior halves.
\end{itemize}

\subsection{Anatomical positional terms}

\begin{itemize}
    \item Superior
    \item Inferior 
    \item Cephalic (towards the head)\footnote{The terms \term{cephalic} and \term{caudal} are often used in embryology, where the head and tail of the embryo are more distinct than in adults.}
    \item Caudal (towards the tail)\footnotemark[\value{footnote}]
    \item Proximal (closer to the point of attachment)
    \item Distal (further from the point of attachment)
    \item Medial
    \item Lateral
    \item Posterior/dorsal
    \item Anterior/ventral
\end{itemize}

\subsection{Anatomical movement terms}

\begin{itemize}
    \item Supination: rotation of the forearm so that the palm faces anteriorly (or superiorly when the arm is outstretched)
    \item Pronation: rotation of the forearm so that the palm faces posteriorly (or inferiorly when the arm is outstretched)
    \item Inversion: rotation of the foot so that the sole faces medially
    \item Eversion: rotation of the foot so that the sole faces laterally
    \item Flexion: movement that decreases the angle between two body parts\footnote{\label{fn:angle} the `angle' being referred to here is the angle in the `natural' movement (e.g. at the hip it is the raising of the leg that decreases the angle on the anterior side, which is the side that is usually considered when describing flexion and extension).}
    \item Extension: movement that increases the angle between two body parts\footnotemark[\value{footnote}]
    \item Abduction: movement away from the midline of the body
    \item Adduction: movement towards the midline of the body
    \item Circumduction: circular movement of a body part, combining flexion, extension, abduction, and adduction
    \item Rotation: movement around a longitudinal axis, either medially (internal rotation) or laterally (external rotation)
    \item Elevation: movement of a body part superiorly
    \item Depression: movement of a body part inferiorly
    \item Protraction: movement of a body part anteriorly
    \item Retraction: movement of a body part posteriorly
    \item Opposition: movement of the thumb to touch the tips of the other fingers
    \item Reposition: movement of the thumb back to its anatomical position
    \item Dorsiflexion: movement of the foot so that the toes point superiorly
    \item Plantar flexion: movement of the foot so that the toes point inferiorly
    \item Palmar flexion: movement of the hand so that the palm faces inferiorly
\end{itemize}

% \input{chapters/physiology}
%
% \part{Systems}
% \input{chapters/cardiovascular}
% \input{chapters/respiratory}

% -------------------------------------------------
% Optional back matter
% -------------------------------------------------

% \backmatter
% \input{appendices/normal-ranges}
% \input{appendices/abbreviations}

\end{document}
